%=============================================================================
% PART I: TRIPLE EQUIVALENCE AND THERMODYNAMICS
%=============================================================================

\subsection{Triple Equivalence and Statistical Mechanics}
\label{part:thermodynamics}

\section{Introduction: The Structure of Bounded Systems}
\label{sec:introduction_bounded}

\input{c:/Users/kunda/Documents/physics/maxwell/conjecture/publication/ideal-gas/sections/bounded-systems.tex}

\section{Categorical Entropy}
\label{sec:categorical_entropy}

\input{c:/Users/kunda/Documents/physics/maxwell/conjecture/publication/ideal-gas/sections/categorical-entropy.tex}

\section{Oscillatory Entropy}
\label{sec:oscillatory_entropy}

\input{c:/Users/kunda/Documents/physics/maxwell/conjecture/publication/ideal-gas/sections/oscillatory-entropy.tex}

\section{Partition Entropy}
\label{sec:partition_entropy}

\input{c:/Users/kunda/Documents/physics/maxwell/conjecture/publication/ideal-gas/sections/partition-entropy.tex}

\section{Enthalpy and Entropy Equivalence}
\label{sec:enthalpy}

\input{c:/Users/kunda/Documents/physics/maxwell/conjecture/publication/ideal-gas/sections/enthalpy.tex}

\section{Temperature from Three Perspectives}
\label{sec:temperature}

\input{c:/Users/kunda/Documents/physics/maxwell/conjecture/publication/ideal-gas/sections/temperature.tex}

\section{Pressure from Three Perspectives}
\label{sec:pressure}

\input{c:/Users/kunda/Documents/physics/maxwell/conjecture/publication/ideal-gas/sections/pressure.tex}

\section{Internal Energy}
\label{sec:internal_energy}

\input{c:/Users/kunda/Documents/physics/maxwell/conjecture/publication/ideal-gas/sections/internal-energy.tex}

\section{Derivation of the Ideal Gas Law}
\label{sec:ideal_gas_law}

\input{c:/Users/kunda/Documents/physics/maxwell/conjecture/publication/ideal-gas/sections/ideal-gas-law.tex}

\section{Maxwell-Boltzmann Distribution}
\label{sec:maxwell}

\input{c:/Users/kunda/Documents/physics/maxwell/conjecture/publication/ideal-gas/sections/maxwell-distribution.tex}

%=============================================================================
% PART II: POINCARE COMPUTING ARCHITECTURE
%=============================================================================

\subsection{Poincar\'e Computing Architecture}
\label{part:computing}

\section{The S-Entropy Coordinate Space}
\label{sec:finite_space}

\input{c:/Users/kunda/Documents/physics/maxwell/conjecture/publication/ideal-gas/sections/finite-space.tex}

\section{Initial States and Computational Problems}
\label{sec:initial_state}

\input{c:/Users/kunda/Documents/physics/maxwell/conjecture/publication/ideal-gas/sections/initial-state.tex}

\section{The Virtual Instrument}
\label{sec:virtual_instrument}

\input{c:/Users/kunda/Documents/physics/maxwell/conjecture/publication/ideal-gas/sections/virtual-instrument.tex}

\section{Categorical Dynamics}
\label{sec:categorical_dynamics}

\input{c:/Users/kunda/Documents/physics/maxwell/conjecture/publication/ideal-gas/sections/categorical-dynamics.tex}

\section{Solution Trajectories and Recurrence}
\label{sec:solution_trajectory}

\input{c:/Users/kunda/Documents/physics/maxwell/conjecture/publication/ideal-gas/sections/solution-trajectory.tex}

\section{Identity Unification}
\label{sec:identity_unification}

\input{c:/Users/kunda/Documents/physics/maxwell/conjecture/publication/ideal-gas/sections/identity-unification.tex}

\section{Completeness and Answer Equivalence}
\label{sec:completeness}

\input{c:/Users/kunda/Documents/physics/maxwell/conjecture/publication/ideal-gas/sections/poincares-machine-completeness.tex}

\section{Categorical Complexity Theory}
\label{sec:complexity}

\input{c:/Users/kunda/Documents/physics/maxwell/conjecture/publication/ideal-gas/sections/complexity.tex}

\section{Exhaustive Computing and Emergent Memory}
\label{sec:exhaustive}

\input{c:/Users/kunda/Documents/physics/maxwell/conjecture/publication/ideal-gas/sections/exhaustive-computing.tex}

\section{Saint-Entropy Thermodynamics}
\label{sec:stellas}

\input{c:/Users/kunda/Documents/physics/maxwell/conjecture/publication/ideal-gas/sections/st-stellas-thermodynamics.tex}

\section{The Categorical Compiler}
\label{sec:compiler}

\input{c:/Users/kunda/Documents/physics/maxwell/conjecture/publication/ideal-gas/sections/categorical-compiler.tex}

%=============================================================================
% PART III: VALIDATION AND APPLICATIONS
%=============================================================================

\subsection{Validation and Applications}
\label{part:validation}

\section{Trajectory Completion in Phase Space}
\label{sec:trajectory}

\input{c:/Users/kunda/Documents/physics/maxwell/conjecture/publication/ideal-gas/sections/trajectory-completion.tex}

\section{Categorical Memory Implementation}
\label{sec:categorical_memory}

\input{c:/Users/kunda/Documents/physics/maxwell/conjecture/publication/ideal-gas/sections/categorical-memory.tex}

\section{Ternary Representation}
\label{sec:ternary}

\input{c:/Users/kunda/Documents/physics/maxwell/conjecture/publication/ideal-gas/sections/ternary-representation.tex}

%=============================================================================
% DISCUSSION AND CONCLUSION
%=============================================================================

\section{Discussion}
\label{sec:discussion}

\subsection{Resolution of Classical Paradoxes}

The triple equivalence framework resolves several long-standing conceptual difficulties in statistical mechanics.

\subsubsection{Resolution-Dependence of Temperature}

Classical kinetic theory defines temperature through the mean square velocity: $T = m\langle v^2\rangle/(3k_B)$. This definition makes temperature dependent on how velocity is measured---the resolution and averaging procedure introduce apparent ambiguity. In the categorical framework, temperature is instead the rate of categorical actualization:
\begin{equation}
T = \frac{\hbar}{k_B}\frac{dM}{dt}
\end{equation}

Categories are discrete and countable; no resolution ambiguity arises. The classical definition emerges as a projection onto the velocity observable, which introduces apparent resolution-dependence through the measurement process. The categorical rate $dM/dt$ is intrinsic to the system and independent of how we choose to observe it.

\subsubsection{Localization of Pressure}

Classical kinetic theory derives pressure from molecular collisions with container walls, suggesting pressure is fundamentally a boundary phenomenon. Yet we routinely measure pressure in the bulk of fluids, far from any boundaries. In the categorical framework, pressure is categorical density:
\begin{equation}
P = k_B T \left(\frac{\partial M}{\partial V}\right)_S
\end{equation}

This is an intrinsic property existing throughout the volume, not localized at boundaries. Wall collisions are one manifestation of categorical density---the rate at which trajectories encounter boundaries---but not its definition. Pressure exists in the bulk because categorical structure exists in the bulk.

\subsubsection{Infinite Velocity Tail}

The Maxwell-Boltzmann distribution $f(v) \propto v^2 e^{-mv^2/(2k_BT)}$ extends to $v \to \infty$, assigning non-zero probability to arbitrarily high velocities. In the categorical framework, the distribution is over discrete categories $m = 0, 1, \ldots, M_{\max}$, where $M_{\max}$ corresponds to $v_{\max} = c$. The distribution is intrinsically bounded:
\begin{equation}
f(m) = \frac{e^{-\beta E_m}}{\sum_{m=0}^{M_{\max}} e^{-\beta E_m}}
\end{equation}

No particle can occupy category $m > M_{\max}$, automatically enforcing $v \leq c$.

\subsection{Unification of Thermodynamics and Computing}

The central insight of this work is that thermodynamics and computation are two perspectives on the same underlying categorical structure. The ideal gas law $PV = Nk_BT$ is simultaneously:
\begin{itemize}
    \item A thermodynamic relation between macroscopic observables
    \item A categorical balance condition in bounded phase space
    \item A constraint on valid solution trajectories in Poincar\'e Computing
\end{itemize}

The processor-oscillator duality ($R_{\text{compute}} = \omega/(2\pi)$) establishes that every oscillator is simultaneously a processor, and vice versa. Hardware timing jitter is not noise to be minimized but categorical information to be exploited.

\subsection{The $\epsilon$-Boundary and Asymptotic Solutions}

Solutions in Poincar\'e Computing are recognized at the $\epsilon$-boundary---exactly one categorical step from closure. This is not approximation but the fundamental nature of categorical computation: categorical irreversibility forbids exact return to the initial state. The ``solution'' IS being one step away; this represents the maximum possible knowledge given G\"odelian constraints.

\section{Conclusion}
\label{sec:conclusion}

We have established that oscillation, categorical distinction, and partition operation are three equivalent descriptions of bounded dynamical systems. From this triple equivalence, we derived:

\textbf{Thermodynamic Results:}
\begin{enumerate}
    \item Entropy in three equivalent forms: $S_{\text{cat}} = k_B M \ln n$, $S_{\text{osc}} = k_B \sum_i \ln(A_i/A_0)$, $S_{\text{part}} = k_B \sum_a \ln(1/s_a)$
    \item Temperature as categorical rate: $T = \hbar(dM/dt)/k_B$
    \item Pressure as categorical density: $P = k_BT(\partial M/\partial V)_S$
    \item The ideal gas law as categorical balance: $PV = Nk_BT$
    \item Maxwell-Boltzmann distribution bounded at $v = c$
\end{enumerate}

\textbf{Computational Results:}
\begin{enumerate}
    \item S-entropy space $\Sspace = [0,1]^3$ satisfies Poincar\'e recurrence
    \item Solutions correspond to recurrent trajectories satisfying constraints
    \item Identity unification: address = processor state = semantic content
    \item Categorical incomparability with Turing computation
    \item Complexity measured in Poincar\'es (categorical completions)
    \item Non-halting dynamics with emergent memory
    \item $3^k$ recursive self-similarity with scale ambiguity
    \item Processor-oscillator duality and processor-memory unification
    \item $\epsilon$-boundary solution recognition
\end{enumerate}

The deeper implication is that thermodynamics and computation are unified through categorical structure in bounded phase space. Temperature measures categorical rate; entropy counts categories; pressure is categorical density; computation is trajectory completion. These are not separate phenomena but projections of a single mathematical structure onto different observables.

All results derive from a single foundational premise: physical systems occupy finite domains. From this boundedness, the triple equivalence structure emerges necessarily, and from the triple equivalence, both statistical mechanics and Poincar\'e Computing follow as mathematical consequences.

% (bibliography managed centrally by monogram.tex)
